\documentclass[tikz,border=10pt]{standalone}
\usetikzlibrary{arrows.meta, positioning, shapes.geometric, shapes.misc, fit, backgrounds, calc, matrix}

\tikzset{
    >=Latex,
    font=\sffamily\large,
    node distance=1.5cm and 3cm,
    dag_result/.style={
        rectangle, 
        draw=green!70!black, 
        fill=green!15, 
        very thick, 
        minimum width=4.5cm, 
        minimum height=1.5cm, 
        text width=5cm,
        align=center
    },
    dag_lemma/.style={
        rounded rectangle, 
        draw=orange!80!black, 
        fill=yellow!20, 
        very thick, 
        minimum width=4.5cm, 
        minimum height=1.5cm, 
        text width=5cm,
        align=center
    },
    dag_model/.style={
        ellipse, 
        draw=blue!70!black, 
        fill=blue!12, 
        very thick, 
        minimum width=4.5cm, 
        minimum height=1.5cm, 
        text width=4.5cm,
        align=center
    },
    dag_unformalized/.style={
        rectangle, 
        draw=gray!60!black, 
        fill=gray!10, 
        very thick, 
        dashed, 
        minimum width=4.5cm, 
        minimum height=1.5cm, 
        text width=5cm,
        align=center
    },
    dag_conditional/.style={
        rectangle, 
        rounded corners,
        draw=orange!80!black, 
        fill=orange!10, 
        very thick, 
        minimum width=4.5cm, 
        minimum height=1.5cm, 
        text width=5cm,
        align=center
    },
    dag_partial/.style={
        rectangle,
        rounded corners,
        draw=yellow!55!black,
        fill=yellow!10,
        very thick,
        dashed,
        minimum width=4.5cm,
        minimum height=1.5cm,
        text width=5cm,
        align=center
    },
    dag_scaffold/.style={
        rectangle,
        draw=gray!55!black,
        fill=white,
        very thick,
        dotted,
        minimum width=4.5cm,
        minimum height=1.5cm,
        text width=5cm,
        align=center
    },
    dag_caveat/.style={
        diamond, 
        draw=red!70!black, 
        fill=red!10, 
        minimum width=4cm, 
        align=center, 
        inner sep=0pt, 
        font=\sffamily\small
    },
    dag_arrow/.style={
        -{Stealth[scale=1.2]}, 
        very thick,
        shorten >=1.5pt,
        shorten <=1.5pt
    },
    dag_dashed_arrow/.style={
        -{Stealth[scale=1.2]}, 
        very thick, 
        dashed,
        shorten >=1.5pt,
        shorten <=1.5pt
    },
    dag_template_legend/.style={
        minimum width=2.6cm,
        minimum height=0.8cm,
        text width=2.45cm,
        inner sep=3pt,
        font=\sffamily\small,
        align=center
    },
    dag_template_note/.style={
        rectangle,
        draw=gray!50,
        fill=gray!5,
        rounded corners,
        text width=13cm,
        align=left,
        font=\sffamily\small,
        inner sep=6pt
    }
}

% Legend Helper
\newcommand{\daglegend}[2]{
    \node[draw, rounded corners, very thick, fit=#1, inner sep=12pt, label=above:\textbf{\Large Legend}] (#2) {};
}


\begin{document}

% Agent visual validation loop:
% 1. Compile this file to PDF after every substantive DAG edit.
% 2. Inspect the PDF or a rendered PNG, not just the TeX source.
% 3. Increase spacing or reroute arrows until no box, legend, note, edge, or label overlaps.
%
% Intended lifecycle:
% - During paper intake, replace the scaffold below with every named result in the
%   paper and the dependency arrows between them.
% - After intake, treat this as a stable topology: update node styles/status/text
%   as proofs close, but avoid adding boxes or arrows unless the original named
%   result inventory or dependency map was actually incomplete.
%
% Layout reminder:
% Use `\dagPaperMetadata`, `\dagPaperLegendRightOfMetadata`, and
% `\begin{dagPaperBody}` so metadata stays top-left, the legend sits to its
% right, and the graph starts below them from the diagram left edge.
%
% Arrow direction rule:
% Draw edges from prerequisite/source node to dependent/target node.  The
% `dag_arrow` and `dag_dashed_arrow` styles use an explicit end-arrow (`->`),
% so the arrowhead should always land on the node that consumes the dependency.
% Do not reverse the coordinate order just to make routing easier; use anchors
% or an orthogonal route through empty space instead.

\begin{tikzpicture}[
  x=1cm,
  y=1cm,
  every node/.append style={outer sep=4pt}
]
% --- Legend: README status vocabulary plus formalized node-type styles. ---
\dagPaperMetadata{[Paper Title]}{[Authors]}{[Publication Venue]}{[Year]}{[Formalized PDF link]}
\dagPaperLegendRightOfMetadata{
\node[dag_result, dag_template_legend] (legRes) at (0,0) {formalized\\result};
\node[dag_lemma, dag_template_legend] (legLem) at (4.2,0) {formalized\\lemma};
\node[dag_model, dag_template_legend] (legDef) at (8.4,0) {formalized\\definition};
\node[dag_caveat_legend] (legCav) at (12.6,0) {formalized\\with caveat};
\node[dag_partial, dag_template_legend] (legPart) at (0,-2.6) {partially\\formalized};
\node[dag_conditional, dag_template_legend] (legCond) at (4.2,-2.6) {conditional};
\node[dag_scaffold, dag_template_legend] (legScaf) at (8.4,-2.6) {scaffold};
\node[dag_unformalized, dag_template_legend] (legNot) at (12.6,-2.6) {not started /\\not formalized};

}
\daglegend{(legRes)(legLem)(legDef)(legCav)(legPart)(legCond)(legScaf)(legNot)}{Legend}

\begin{dagPaperBody}

% --- Layout guidance for future agents. ---
\node[dag_template_note] (Guide) {
  Replace these scaffold nodes with the paper's named definitions, lemmas,
  propositions, theorems, and corollaries. Keep nodes on the grid below:
  columns are spaced by 6.4cm and rows by 3.2cm, which is wide enough for the
  default 5cm node text widths. Prefer straight or orthogonal arrows (`--`,
  `|-`, `-|`) and route through empty grid cells when a dependency skips a row.
  Draw from prerequisite to dependent result so the arrowhead points at the
  consumer.  For short horizontal gaps, use explicit anchors like `(A.east) --
  (B.west)` and increase spacing if the arrowhead visually merges with a node.
  After every layout change, render and inspect the PDF; if anything overlaps,
  increase row or column spacing before adding complicated curves. Once the
  initial named-result map is complete, prefer README-status style updates over
  adding new boxes or arrows.
};

% --- Non-overlapping proof graph scaffold. ---
% Status update pattern:
% - Change the leading style only, e.g. dag_unformalized -> dag_conditional ->
%   dag_lemma/dag_result, and update the short text if the remaining assumption changed.
% - Keep node names and arrows stable after the initial intake map is complete.
%
% Node/edge status semantics:
% - Green nodes (`dag_result`) mean the displayed theorem/result is fully
%   formalized under its stated Lean assumptions. A green theorem may bypass an
%   incomplete paper-route lemma, but then that incomplete lemma must not be
%   shown as a required solid input.
% - Solid arrows (`dag_arrow`) are verified dependencies in the formalized proof
%   path.
% - Dashed arrows (`dag_dashed_arrow`) are for unresolved, conditional, caveat,
%   paper-roadmap, or bypassed-paper-route dependencies. If a dashed edge points
%   into a green node, document in the node/README that the edge is not required
%   by the formal proof; otherwise the target should be conditional/caveated.
\node[dag_model] (Model) at (0,-10.2) {
  \textbf{Definitions / Model} \\
  Source primitives
};
\node[dag_lemma] (LemmaA) at (6.4,-10.2) {
  \textbf{Lemma A} \\
  First reusable step
};
\node[dag_lemma] (LemmaB) at (12.8,-10.2) {
  \textbf{Lemma B} \\
  Second reusable step
};

\node[dag_conditional] (Bridge) at (6.4,-13.4) {
  \textbf{Bridge Theorem} \\
  Conditional paper-facing reduction
};
\node[dag_unformalized] (Open) at (12.8,-13.4) {
  \textbf{Open Source Result} \\
  Name exact remaining gap
};

\node[dag_result] (Main) at (6.4,-16.6) {
  \textbf{Main Theorem} \\
  Closed paper-facing endpoint
};

\draw[dag_arrow] (Model) -- (LemmaA);
\draw[dag_arrow] (LemmaA) -- (LemmaB);
\draw[dag_arrow] (LemmaA) -- (Bridge);
\draw[dag_dashed_arrow] (LemmaB) -- (Open);
\draw[dag_arrow] (Bridge) -- (Main);
\draw[dag_dashed_arrow] (Open) |- (Main);
\end{dagPaperBody}
\end{tikzpicture}
\end{document}
